%%%%%%
%
% $Author: Abishek $
% $Date: 12.01.2025 $
% $Path: ML24-02-Gait-Pattern-Recognition\report\Contents\en$
% $Version: 1.0 $
%
%
% !TeX encoding = utf8
% !TeX root = Rename
% !TeX TXS-program:bibliography = txs:///bibtex
%
%%%%%%

\section{Given Documentation}

\subsection{Use of}
The provided documentation served as a foundational reference for the hardware aspects of the project. It included technical details, descriptions, and code snippets relevant to the implementation and analysis of the hardware components. Key details about hardware like the Arduino Nano 33 BLE Sense and sensors, including the LSM9DS1 accelerometer, were integrated directly into the report. For instance:
\begin{itemize}
	\item The technical specifications table for the Arduino Nano 33 BLE Sense was reproduced to highlight its features, such as operating voltage, power consumption, and sensor integration.
	\item Sensor descriptions, such as the accelerometer's ability to measure acceleration across three axes (X, Y, Z), were adopted to ensure consistency in terminology and units.
\end{itemize}

\subsection{Same Notation}
To ensure uniformity and clarity, the same notations and units used in the original reference document were followed throughout the report. Examples include:
\begin{itemize}
	\item Units such as ±2g, ±4g for acceleration ranges and terms like “g” for gravitational force were consistently applied in both tabular and textual formats.
	\item Technical notations for current consumption (mA, µA), resolution (bits), and temperature range (°C) were maintained without alteration.
	\item Diagrams and illustrations, such as the accelerometer’s X, Y, Z axes, used the same labeling and directionality, recreated using TikZ for consistency.
\end{itemize}

\subsection{Supplements}
Visual aids such as diagrams, figures, and tables were included to support the textual descriptions and provide better insights into the hardware functionality. Examples include:
\begin{itemize}
	\item The accelerometer axes diagram and the specifications table for the LSM9DS1 were directly included, recreated using LaTeX tools for a uniform style.
	\item Figures such as the Arduino Nano 33 BLE Sense component overview and pin configuration diagrams were integrated to provide readers with a clear understanding of hardware layout and usage.
	\item Code snippets for calibration and real-time data collection were incorporated to demonstrate practical implementation, maintaining the same structure and functionality as the reference.
\end{itemize}

\subsection{Same Structure}
The structure of the documentation was mirrored to maintain a logical flow and make the report easier to navigate. Examples include:
\begin{itemize}
	\item Sections such as \textit{IMU Overview}, \textit{Sensor Specifications}, and \textit{Library Installation} were arranged in the same hierarchy as the reference document.
	\item Subsections for each sensor (e.g., \textit{LSM9DS1, HTS221, MP34DT05}) followed the same sequence, ensuring a systematic approach.
	\item Code snippets and examples, such as IMU calibration and tilt detection, were positioned in the corresponding sections for consistency.
\end{itemize}

\subsection*{Examples of Alignment}

\textbf{Direct Usage:}
\begin{itemize}
	\item The specifications table for the LSM9DS1 accelerometer was reproduced with no changes.
	\item Descriptions of power consumption modes (Normal, Low-Power, Standby) were directly used.
	\item Diagrams illustrating the axes of the accelerometer and their functionality were adopted with consistent notation.
\end{itemize}

\textbf{Integrated Adaptations:}
\begin{itemize}
	\item Hardware descriptions were contextualized for the report’s specific focus, such as applications in motion detection and tilt sensing.
	\item Explanations of sensor calibration were expanded to include additional use cases and step-by-step implementation details.
\end{itemize}

\textbf{Custom Extensions:}
\begin{itemize}
	\item Sections on implementing machine learning applications with the Arduino Nano 33 BLE Sense were introduced.
	\item Code examples for tilt angle detection and real-time Serial Monitor outputs were customized to highlight the project’s application-specific contributions.
\end{itemize}
\section{Improvements}

During the project development, we continuously improved our work based on iterative refinements and feedback. These modifications significantly enhanced the accuracy, clarity, and usability of our documentation, technical implementation, and the overall project. Below, we outline the key improvements made, their justification, and the resulting benefits.

\subsection{Minor Changes}
Several small adjustments were made to refine the documentation, improve readability, and enhance usability:
\begin{itemize}
	\item Divided sensor types (accelerometers, gyroscopes, and magnetometers) into bullet points for better readability.
	\item Enhanced descriptions of individual sensor contributions, explicitly detailing their roles.
	\item Standardized chart colors to improve visual clarity and logical grouping of data points.
	\item Reformatted and added headers to the code for better readability and navigation.
	\item Included academic and technical citations to justify challenges, design decisions, and solutions.
	\item Simplified the classification task from three features (normal, Parkinson’s, and stroke) to two (normal vs. Parkinson’s) to focus on the project scope.
	\item Integrated a troubleshooting guide and step-by-step resolutions for common issues, minimizing downtime.
\end{itemize}

\subsection{List of Minor Changes}
\begin{enumerate}
	\item Reformatted sensor descriptions into bullet points.
	\item Added references to support challenges and solutions.
	\item Improved the material list with real images and direct purchase links for easier sourcing.
	\item Added comments and structured headers to the source code.
	\item Updated inconsistent charts with standardized colors for logical data grouping.
\end{enumerate}

\subsection{Major Improvements}
The project underwent significant improvements in the following areas:

\subsubsection{Hardware Adaptation}
The documentation initially focused on the Arduino Nano 33 BLE Sense, but the project required the Lite version. The hardware documentation was adapted to reflect the capabilities and limitations of the Lite version:
\begin{itemize}
	\item Differences between the standard and Lite versions were highlighted, such as the absence of certain sensors.
	\item Adjustments were made to ensure compatibility with the Lite version while retaining functionality.
\end{itemize}

\subsubsection{Machine Learning Pipeline}
Initially, Random Forest was planned as the core algorithm, but Convolutional Neural Networks (CNNs) were adopted for better performance. This change improved:
\begin{itemize}
	\item Data flow and feature extraction capabilities.
	\item Accuracy in classifying features (normal vs. Parkinson’s).
	\item Clarity in explaining the machine learning pipeline, including detailed descriptions of training, testing, and inference processes.
\end{itemize}

\subsubsection{Maintenance Monitoring Upgrade}
The original monitoring system required offline maintenance. This was transitioned to an online remote maintenance system, offering:
\begin{itemize}
	\item Real-time monitoring capabilities via cloud-based solutions.
	\item Improved user accessibility and reduced dependency on physical visits.
\end{itemize}

\subsubsection{Documentation Enhancements}
\begin{itemize}
	\item Improved visual representation of hardware components using AutoCAD prototypes.
	\item Created detailed pin descriptions with TikZ diagrams for better understanding of wiring and connections.
	\item Developed a battery connection guide to avoid wiring mistakes, ensuring hardware safety.
	\item Addressed power management concerns by providing a schematic and explanation for using USB and battery power simultaneously.
	\item Added a detailed electrical circuit explanation, including a component description table for easier troubleshooting and assembly.
\end{itemize}

\subsection{Impact and Benefits of These Changes}
These improvements significantly enhanced the project by:
\begin{itemize}
	\item \textbf{Improved Clarity}: Enhanced diagrams, structured explanations, and refined documentation reduced ambiguity.
	\item \textbf{Stronger Justification}: Academic citations validated the design choices and solutions.
	\item \textbf{Better Usability}: A refined user manual, troubleshooting guide, and material list improved ease of use and implementation.
	\item \textbf{Enhanced Technical Accuracy}: Detailed proportional circuit diagrams and ML pipeline documentation provided precise implementation guidelines.
	\item \textbf{Reduced Errors}: Visual guides, structured explanations, and step-by-step troubleshooting minimized assembly and deployment errors.
\end{itemize}

By addressing these points, the project evolved into a more robust, user-friendly, and professionally documented system, adhering to best practices in hardware integration, machine learning, and engineering.

